\chapter{Resultados esperados y gestión de riesgos}

\section{Resultados esperados}

El proyecto espera producir: (i) una caracterización documentada de conjuntos de imágenes dermatológicas y sus límites de representatividad; (ii) un clasificador que use representaciones aprendidas mediante SSL y un clasificador lineal posterior claramente descrito; (iii) una evaluación reproducible frente a controles y, cuando se complete el protocolo comparable, una línea base supervisada; y (iv) un prototipo experimental de apoyo a la revisión de imágenes.

No se fijan umbrales de desempeño como resultados garantizados. Las métricas se reportarán con su protocolo, incertidumbre y límites de dominio. El sistema no sustituye al profesional ni representa una herramienta clínica validada.

\section{Riesgos y aseguramiento de calidad}

\begin{longtable}{p{0.24\textwidth}p{0.34\textwidth}p{0.34\textwidth}}
\caption{Riesgos prioritarios y medidas de control.}\label{tab:riesgos-principales}\\
\toprule
Riesgo & Consecuencia para la evidencia & Control previsto \\
\midrule
\endfirsthead
\toprule
Riesgo & Consecuencia para la evidencia & Control previsto \\
\midrule
\endhead
Solapamiento de imágenes de una lesión entre particiones & Fuga de información y métricas optimistas & Separar por \texttt{lesion\_id}; congelar listas y hashes antes de entrenar. \\
\addlinespace
Bootstrap por imagen en HAM10000 & Intervalos demasiado estrechos por correlación intralesión & Reportar la limitación y hacer sensibilidad agrupada por lesión antes de afirmaciones confirmatorias. \\
Conjuntos de datos y modalidades diferentes & Transferencia aparente o deterioro fuera de dominio & Describir modalidad, fuente y población; reportar PAD-UFES-20 por separado. \\
Línea base supervisada con partición distinta & Comparación no pareada y no atribuible solo al régimen de aprendizaje & Igualar conjunto de datos, arquitectura, partición, presupuesto y métricas antes de concluir. \\
ITA presentado como fototipo & Conclusiones inválidas sobre equidad por FST & Nombrar ITA correctamente; no afirmar FST sin anotación experta y validación. \\
Procedencia incompleta de corridas/pesos & Dificultad para reproducir o auditar resultados & Registrar revisión, configuración, semillas, hashes de particiones/pesos y estado W\&B. \\
\bottomrule
\end{longtable}

\section{Ética, privacidad y límites de uso}

El proyecto usa conjuntos públicos, sujetos a sus respectivas licencias, y no recoge imágenes de pacientes. Antes de reutilizar o redistribuir imágenes, se comprobarán los términos de cada conjunto. El prototipo se plantea como demostración académica en entorno de laboratorio; cualquier recolección de datos clínicos, prueba con pacientes o despliegue requiere revisión institucional y un protocolo distinto.

\noindent\textit{[Pendiente de confirmación con dirección y la instancia ética de la UNAB: requisitos institucionales para investigación con conjuntos públicos y condiciones de uso o redistribución de cada conjunto.]}
